
\FigureLayoutDeclare{img/2022/shanghai_spring/q11_solution_fig1.png}{8.0cm}{0bp 0bp 0bp 0bp}

\chapter{2022年上海卷（春）}

\section{填空题}

\begin{problem}
已知 \(z=2+i\)，则 \(\bar z=\fillinblank{}\).
\end{problem}

\begin{problem}
已知 \(A=(-1,2)\)，\(B=(1,3)\)，则 \(A\cap B=\fillinblank{}\).
\end{problem}

\begin{problem}
不等式 \(\dfrac{x-1}{x}<0\) 的解集为\fillinblank{}.
\end{problem}

\begin{problem}
已知 \(\tan\alpha=3\)，则 \(\tan\left(\alpha+\dfrac\pi4\right)=\fillinblank{}\).
\end{problem}

\begin{problem}
已知方程组 \(\begin{cases}
x+my=2,\\
mx+16y=8
\end{cases}\) 有无穷解，则 \(m\) 的值为\fillinblank{}.
\end{problem}

\begin{problem}
已知函数 \(f(x)=x^3\) 的反函数为 \(y=f^{-1}(x)\)，则 \(f^{-1}(27)=\fillinblank{}\).
\end{problem}

\begin{problem}
在 \(\left(x^3+\dfrac1x\right)^{12}\) 的展开式中，含 \(\dfrac1{x^4}\) 项的系数为\fillinblank{}.
\end{problem}

\begin{problem}
在 \(\triangle ABC\) 中，\(\angle A=\dfrac\pi3\)，\(AB=2\)，\(AC=3\)，则 \(\triangle ABC\) 的外接圆半径为\fillinblank{}.
\end{problem}

\begin{problem}
已知有 \(1\)，\(2\)，\(3\)，\(4\) 四个数字组成无重复数字，则比 \(2134\) 大的四位数的个数为\fillinblank{}.
\end{problem}

\begin{problem}
在 \(\triangle ABC\) 中，\(\angle C=\dfrac\pi2\)，\(AC=BC=2\)，\(M\) 为 \(AC\) 的中点，\(P\) 在线段 \(AB\) 上，则 \(\overrightarrow{MP}\cdot\overrightarrow{CP}\) 的最小值为\fillinblank{}.
\end{problem}

\begin{problem}
已知双曲线 \(\dfrac{x^2}{a^2}-y^2=1\)，其中 \(a>0\)，双曲线右支上有任意两点 \(P_1(x_1,y_1)\)，\(P_2(x_2,y_2)\)，满足 \(x_1x_2-y_1y_2>0\) 恒成立，则 \(a\) 的取值范围是\fillinblank{}.
\end{problem}

\begin{problem}
已知 \(f(x)\) 为奇函数，当 \(x\in(0,1]\) 时，\(f(x)=\ln x\)，且 \(f(x)\) 关于直线 \(x=1\) 对称. 设 \(f(x)=x+1\) 的正数解依次为 \(x_1\)，\(x_2\)，\(x_3\)，\(\ldots\)，\(x_n\)，\(\ldots\)，则 \(\displaystyle\lim_{n\to\infty}(x_{n+1}-x_n)=\fillinblank{}\).
\end{problem}

\section{单选题}

\begin{problem}
下列函数定义域为 \(\R\) 的是（\quad）

\choices
    {\(y=x^{-\frac12}\)}
    {\(y=x^{-1}\)}
    {\(y=x^{\frac13}\)}
    {\(y=x^{\frac12}\)}
\end{problem}

\begin{problem}
已知 \(a>b>c>d\)，下列选项中正确的是（\quad）

\choices
    {\(a+d>b+c\)}
    {\(a+c>b+d\)}
    {\(ad>bc\)}
    {\(ac>bd\)}
\end{problem}

\begin{problem}
如图，上海海关大楼的上面可以看作一个正四棱柱，四个侧面有四个时钟，则相邻两个时钟的时针从 \(0\) 时转到 \(12\) 时（含 \(0\) 时不含 \(12\) 时）的过程中，能够相互垂直（\quad）次.

\bitmapfigure[max width=0.22\linewidth]{img/2022/shanghai_spring/q15_fig1.png}

\choices
    {\(0\)}
    {\(2\)}
    {\(4\)}
    {\(12\)}
\end{problem}

\begin{problem}
已知等比数列 \(\{a_n\}\) 的前 \(n\) 项和为 \(S_n\)，前 \(n\) 项积为 \(T_n\)，则下列选项判断正确的是（\quad）

\choices
    {若 \(S_{2022}>S_{2021}\)，则数列 \(\{a_n\}\) 是递增数列}
    {若 \(T_{2022}>T_{2021}\)，则数列 \(\{a_n\}\) 是递增数列}
    {若数列 \(\{S_n\}\) 是递增数列，则 \(a_{2022}\ge a_{2021}\)}
    {若数列 \(\{T_n\}\) 是递增数列，则 \(a_{2022}\ge a_{2021}\)}
\end{problem}

\section{解答题}

\begin{problem}
如图，在圆柱 \(OO_1\) 中，底面半径为 \(1\)，\(AA_1\) 为圆柱 \(OO_1\) 的母线.

\begin{enumerate}
\item 若 \(AA_1=4\)，\(M\) 为 \(AA_1\) 的中点，求直线 \(MO_1\) 与底面的夹角大小；
\item 若圆柱的轴截面为正方形，求该圆柱的侧面积和体积.
\end{enumerate}

\bitmapfigure[max width=0.33\linewidth]{img/2022/shanghai_spring/q17_fig1.png}
\end{problem}

\begin{problem}
已知数列 \(\{a_n\}\)，\(a_2=1\)，\(\{a_n\}\) 的前 \(n\) 项和为 \(S_n\).

\begin{enumerate}
\item 若 \(\{a_n\}\) 为等比数列，\(S_2=3\)，求 \(\displaystyle\lim_{n\to\infty}S_n\)；
\item 若 \(\{a_n\}\) 为等差数列，公差为 \(d\)，对任意 \(n\in\N^*\)，均满足 \(S_{2n}\ge n\)，求 \(d\) 的取值范围.
\end{enumerate}
\end{problem}

\begin{problem}
如图，矩形 \(ABCD\) 区域内，\(D\) 处有一棵古树，为保护古树，以 \(D\) 为圆心，\(DA\) 为半径划定圆 \(D\) 作为保护区域，已知 \(AB=30\mathrm{~m}\)，\(AD=15\mathrm{~m}\)，点 \(E\) 为 \(AB\) 上的动点，点 \(F\) 为 \(CD\) 上的动点，满足 \(EF\) 与圆 \(D\) 相切.

\begin{enumerate}
\item 若 \(\angle ADE=20^\circ\)，求 \(EF\) 的长；
\item 当点 \(E\) 在 \(AB\) 的什么位置时，梯形 \(FEBC\) 的面积有最大值，最大面积为多少？
\end{enumerate}

\bitmapfigure[max width=0.42\linewidth]{img/2022/shanghai_spring/q19_fig1.png}

（长度精确到 \(0.1\mathrm{~m}\)，面积精确到 \(0.01\mathrm{~m}^2\)）
\end{problem}

\begin{problem}
在椭圆 \(\Gamma:\dfrac{x^2}{a^2}+y^2=1\) 中，直线 \(l:x=a\) 上有两点 \(C\)，\(D\)（\(C\) 点在第一象限），左顶点为 \(A\)，下顶点为 \(B\)，右焦点为 \(F\).

\begin{enumerate}
\item 若 \(\angle AFB=\dfrac\pi6\)，求椭圆 \(\Gamma\) 的标准方程；
\item 若点 \(C\) 的纵坐标为 \(2\)，点 \(D\) 的纵坐标为 \(1\)，则 \(BC\) 与 \(AD\) 的交点是否在椭圆上？请说明理由；
\item 已知直线 \(BC\) 与椭圆 \(\Gamma\) 相交于点 \(P\)，直线 \(AD\) 与椭圆 \(\Gamma\) 相交于点 \(Q\)，若 \(P\) 与 \(Q\) 关于原点对称，求 \(\lvert CD\rvert\) 的最小值.
\end{enumerate}
\end{problem}

\begin{problem}
已知函数 \(f(x)\)，甲变化：\(f(x)-f(x-t)\)；乙变化：\(\lvert f(x+t)-f(x)\rvert\)，\(t>0\).

\begin{enumerate}
\item 若 \(t=1\)，\(f(x)=2^x\)，\(f(x)\) 经甲变化得到 \(g(x)\)，求方程 \(g(x)=2\) 的解；
\item 若 \(f(x)=x^2\)，\(f(x)\) 经乙变化得到 \(h(x)\)，求不等式 \(h(x)\le f(x)\) 的解集；
\item 若 \(f(x)\) 在 \((-infty,0)\) 上单调递增，将 \(f(x)\) 先进行甲变化得到 \(u(x)\)，再将 \(u(x)\) 进行乙变化得到 \(h_1(x)\)；将 \(f(x)\) 先进行乙变化得到 \(v(x)\)，再将 \(v(x)\) 进行甲变化得到 \(h_2(x)\)，若对任意 \(t>0\)，总存在 \(h_1(x)=h_2(x)\) 成立，求证：\(f(x)\) 在 \(\R\) 上单调递增.
\end{enumerate}
\end{problem}

